\pdfoutput=1
% ================================================================
% SCX — Clinical Trials as Multi-Expert Audit
% ================================================================
% Compile: pdflatex -interaction=nonstopmode main.tex
% ================================================================
\documentclass[12pt,a4paper]{article}

% ================================================================
% Language & Encoding
% ================================================================
\usepackage[english]{babel}
\usepackage[T1]{fontenc}

% ================================================================
% Page Geometry
% ================================================================
\usepackage[a4paper,top=2cm,bottom=2cm,left=2.5cm,right=2.5cm]{geometry}

% ================================================================
% Math Core
% ================================================================
\usepackage{amsmath,amssymb,amsthm}
\usepackage{mathtools}
\usepackage{bm}
\usepackage{bbm}

% ================================================================
% Graphics & Color
% ================================================================
\usepackage{graphicx}
\usepackage{tikz}
\usepackage{xcolor}
\usepackage{float}

% ================================================================
% Tables & Lists
% ================================================================
\usepackage{booktabs}
\usepackage{enumitem}
\usepackage{paralist}

% ================================================================
% Algorithms
% ================================================================
\usepackage{algorithm}
\usepackage{algpseudocode}

% ================================================================
% Hyperlinks & Cross-References
% ================================================================
\usepackage[colorlinks=true, citecolor=blue, urlcolor=blue]{hyperref}
\usepackage{cleveref}

% ================================================================
% Theorem Environments (English)
% ================================================================
\newtheorem{theorem}{Theorem}[section]
\newtheorem{lemma}[theorem]{Lemma}
\newtheorem{proposition}[theorem]{Proposition}
\newtheorem{corollary}[theorem]{Corollary}
\newtheorem{definition}[theorem]{Definition}
\newtheorem{conjecture}[theorem]{Conjecture}
\newtheorem{assumption}[theorem]{Assumption}

\theoremstyle{remark}
\newtheorem{remark}[theorem]{Remark}
\newtheorem{example}[theorem]{Example}
\newtheorem{protocol}[theorem]{Protocol}

% ================================================================
% Math Shortcuts — Blackboard Bold
% ================================================================
\newcommand{\R}{\mathbb{R}}
\newcommand{\C}{\mathbb{C}}
\newcommand{\N}{\mathbb{N}}
\newcommand{\Z}{\mathbb{Z}}
\newcommand{\Q}{\mathbb{Q}}
\newcommand{\E}{\mathbb{E}}
\newcommand{\Pbb}{\mathbb{P}}

% ================================================================
% Math Shortcuts — Calligraphic
% ================================================================
\newcommand{\D}{\mathcal{D}}
\newcommand{\F}{\mathcal{F}}
\newcommand{\G}{\mathcal{G}}
\newcommand{\T}{\mathcal{T}}
\newcommand{\loss}{\mathcal{L}}
\newcommand{\Obs}{\mathcal{O}}
\newcommand{\M}{\mathcal{M}}
\newcommand{\A}{\mathcal{A}}
\newcommand{\B}{\mathcal{B}}
\newcommand{\I}{\mathcal{I}}
\newcommand{\J}{\mathcal{J}}
\newcommand{\X}{\mathcal{X}}
\newcommand{\K}{\mathcal{K}}
\newcommand{\Sg}{\mathcal{S}}
\newcommand{\U}{\mathcal{U}}
\newcommand{\W}{\mathcal{W}}
\newcommand{\Trial}{\mathcal{T}}
\newcommand{\AuditGrp}{\mathfrak{G}_{\text{audit}}}
\newcommand{\ExpertSet}{\EuScript{E}}
\newcommand{\CGC}{\mathcal{C}}

% ================================================================
% SCX Framework Macros
% ================================================================
\newcommand{\SCX}{\textsf{SCX}}
\newcommand{\Yajie}{\textsf{Yajie}}
\newcommand{\Spring}{\textsf{Spring}}
\newcommand{\Cercis}{\textsf{Cercis}}
\newcommand{\Situs}{\textsf{Situs}}
\newcommand{\MoE}{\textsf{MoE}}
\newcommand{\Placebo}{\textsf{Placebo}}
\newcommand{\Meta}{\textsf{Meta}}
\newcommand{\RCT}{\textsf{RCT}}
\newcommand{\EBM}{\textsf{EBM}}

% ================================================================
% Operators
% ================================================================
\DeclareMathOperator{\argmax}{arg\,max}
\DeclareMathOperator{\argmin}{arg\,min}
\DeclareMathOperator{\sgn}{sgn}
\DeclareMathOperator{\diag}{diag}
\DeclareMathOperator{\spec}{spec}
\DeclareMathOperator{\softmax}{softmax}
\DeclareMathOperator{\topk}{top\text{-}k}
\DeclareMathOperator{\Cov}{Cov}
\DeclareMathOperator{\Var}{Var}
\DeclareMathOperator{\Tr}{Tr}
\DeclareMathOperator{\KL}{KL}
\DeclareMathOperator{\rank}{rank}
\DeclareMathOperator{\supp}{supp}
\DeclareMathOperator{\RR}{RR}
\DeclareMathOperator{\OR}{OR}
\DeclareMathOperator{\HR}{HR}
\DeclareMathOperator{\NNT}{NNT}
\DeclareMathOperator{\se}{se}
\DeclareMathOperator{\Ihet}{I^2}
\DeclareMathOperator{\AuditDepth}{AuditDepth}

% ================================================================
% Proof Status Markers
% ================================================================
\newcommand{\rigorFull}{\textbf{[Full Proof]}}
\newcommand{\rigorPartial}{\textbf{[Partial Proof]}}
\newcommand{\rigorSketch}{\textbf{[Proof Sketch]}}
\newcommand{\openproblem}{\textbf{[Open Problem]}}

% ================================================================
% Honest Critique
% ================================================================
\newcommand{\honestcrit}[1]{\textcolor{red}{\textbf{[Honest Critique]} #1}}

% ================================================================
% Document Metadata
% ================================================================
\title{\textbf{Clinical Trials as Multi-Expert Audit:\\
Phase III, Meta-Analysis, and the Cercis Foundation\\
of Evidence-Based Medicine}}
\author{SCX}
\date{\today}

% ================================================================
\begin{document}
\maketitle

\begin{abstract}
Evidence-based medicine (EBM) rests on a hierarchy of evidence: systematic
reviews at the apex, followed by randomized controlled trials (RCTs), cohort
studies, case series, and finally expert opinion at the base. We prove that
this hierarchy is precisely a spectrum of \textbf{audit depth} in the \SCX{}
framework, with clinical trials functioning as principled multi-expert
($M > 1$) audits of treatment efficacy. The core contributions are nine
formal results establishing clinical trials as the original and most
consequential application of multi-expert audit theory. (1) A Phase III
randomized controlled trial is an \textbf{$M$-expert audit} where patients,
investigators, and data monitors collectively assess the treatment effect,
with blinding and randomization serving as gauge-fixing mechanisms.
(2) The placebo control arm is \textbf{gauge zero}: placebo response is a
gauge artifact that arises from patient and observer expectations, and
blinding removes it by making the treatment assignment gauge-invariant.
(3) Meta-analysis computes the \textbf{Cercis score} across independent
trials: the overall effect estimate is the quality-weighted average, and
the heterogeneity statistic $I^2$ measures \textbf{gauge misalignment}
across studies. (4) Diagnostic disagreement among radiologists is a
\textbf{high-Cercis signal}: second opinions correspond to an $M+1$ audit
that exponentially reduces the false-positive rate. (5) The EBM hierarchy
is an \textbf{audit depth function} $\AuditDepth(\text{study type})$ that
maps study design to the effective number of independent expert auditors.
(6) Eminence-based medicine (``because I said so'') is the $M = 1$
catastrophe: a single expert self-certifies without external audit,
analogous to a gatekeeper with infinite concealment. (7) The reproducibility
crisis is an \textbf{audit failure}: $p$-hacking corresponds to gauge
manipulation, and publication bias corresponds to censored audit data.
(8) $p$-hacking shifts the audit gauge, inflating the false-positive rate
by exactly the factor predicted by the \SCX{} gauge transformation theory.
(9) Cercis-based trial design gives the optimal number of experts $M^*$
for a given effect size and desired statistical power. A complete numerical
verification suite is provided in \texttt{verify\_medicine\_audit.py}.
\end{abstract}


% ================================================================
\section{Introduction}
\label{sec:intro}
% ================================================================

Clinical medicine confronts a fundamental epistemic challenge: how does one
know whether a treatment works? The answer, developed over the past seven
decades, is the randomized controlled trial~\cite{hill1965environment,
fisher1935design}. An RCT enrolls patients, randomly assigns them to
treatment or control, blinds participants and investigators to the
assignment, and measures a pre-specified outcome. The statistical test at
the end---did the treatment group do better than chance?---is the definitive
arbiter.

This is an audit problem.

The \SCX{} framework (State-Conditioned eXpertise)~\cite{scx_theory}
provides a formal theory of multi-expert consensus, bias detection, and
information concealment. Its core insight is that the \Cercis{} score
$S(s) = Q(s) + \eta(t) \cdot N(s)$ decomposes any audited state into a
quality (agreement) component and a novelty (distinguishability) component.
An expert with a non-zero gatekeeper parameter conceals their true state
behind an information barrier. In this paper, we show that the entire
edifice of evidence-based medicine---from the Phase III RCT to systematic
review and meta-analysis---is a manifestation of the \SCX{} audit framework,
with clinical trials serving as the paradigmatic $M > 1$ audit.

The central claim is that the evidence hierarchy of EBM is not merely a
convenient ranking of study designs but a mathematically precise
\textbf{audit depth hierarchy}. At the top, systematic reviews and
meta-analyses aggregate information across many independent trials
($M \gg 1$). At the bottom, expert opinion is a single expert's assertion
($M = 1$), which the \SCX{} framework identifies as fundamentally
unaudited---it is eminence-based medicine, not evidence-based medicine.

\medskip
\noindent\textbf{Why this reframing matters.} The mapping from clinical
trials to \SCX{} audit yields quantitative insights:

\begin{enumerate}[leftmargin=*, label=(\arabic*)]
  \item \textbf{Placebo as gauge artifact.} The placebo response is not a
  nuisance to be minimized but a \textbf{gauge artifact}---an expectation
  bias that the blinding procedure (our gauge-fixing transformation)
  removes. This explains why placebo effects are largest in subjective
  outcomes and vanish under objective measurements: subjective outcomes
  are gauge-dependent.
  \item \textbf{Meta-analysis as Cercis computation.} The standard
  meta-analytic fixed-effect estimate $\hat{\theta}_{\text{FE}}$ is exactly
  the quality-weighted Cercis score across trials, with weights inversely
  proportional to the variance (inverse-variance weighting). Random-effects
  meta-analysis corresponds to the \Cercis{} annealing schedule $\eta(t)$.
  \item \textbf{$I^2$ as gauge misalignment.} The heterogeneity statistic
  $I^2$---the proportion of variance due to between-study differences rather
  than chance---is exactly the \textbf{gauge misalignment measure} across
  the expert ensemble.
  \item \textbf{Second opinion as $M+1$ audit.} Diagnostic disagreement
  among radiologists (kappa = 0.4--0.6 in mammography screening) is not
  a failure of medicine but a \textbf{high-Cercis signal}: the disagreement
  rate equals the rate at which a second expert detects gauge-invariant
  information that the first expert missed.
\end{enumerate}

\medskip
\noindent\textbf{Structure of the paper.}
Section~\ref{sec:phase3} establishes Phase III trials as $M>1$ audits.
Section~\ref{sec:placebo} proves the placebo-as-gauge-zero theorem.
Section~\ref{sec:meta} formalizes meta-analysis as \Cercis{} computation.
Section~\ref{sec:diagnostic} analyzes diagnostic disagreement as the
\Cercis{} signal. Section~\ref{sec:hierarchy} maps the EBM evidence
hierarchy to audit depth. Section~\ref{sec:eminence} proves the
$M=1$ catastrophe of eminence-based medicine.
Section~\ref{sec:reproducibility} diagnoses the reproducibility crisis as
audit failure. Section~\ref{sec:trial_design} gives Cercis-based trial
design. Section~\ref{sec:discussion} discusses implications and open
problems.


% ================================================================
\section{Phase III Trials as $M>1$ Audit}
\label{sec:phase3}
% ================================================================

\subsection{The Architecture of the Randomized Controlled Trial}

\begin{definition}[Clinical Trial as Audit Configuration]
\label{def:clinical_trial}
A \textbf{clinical trial audit configuration} is a tuple
$\mathcal{T} = (\mathcal{P}, \mathcal{I}, \mathcal{D}, \mathcal{E})$ where:
\begin{itemize}[leftmargin=*]
  \item $\mathcal{P} = \{p_1, \ldots, p_n\}$ is the set of \textbf{patients},
  each providing an outcome measurement $y_i \in \R$;
  \item $\mathcal{I} = \{i_1, \ldots, i_m\}$ is the set of
  \textbf{investigators}, each providing assessments, measurements, and
  interpretations of the outcomes;
  \item $\mathcal{D} = \{d_1, \ldots, d_k\}$ is the \textbf{data monitoring}
  infrastructure (DSMB, statisticians, auditors) that checks integrity;
  \item $\mathcal{E}$ is the \textbf{evidence function} mapping the trial
  data to a summary statistic $\hat{\theta}$ (e.g., log odds ratio, hazard
  ratio, mean difference) with standard error $\sigma$.
\end{itemize}
The \textbf{expert ensemble} of the trial is the union
$\mathcal{E}_{\text{trial}} = \mathcal{P} \cup \mathcal{I} \cup \mathcal{D}$,
and the \textbf{audit multiplicity} is $M = |\mathcal{E}_{\text{trial}}|$.
\end{definition}

The key insight is that a clinical trial is not a single expert's opinion
but an assembly of many interdependent experts. Each patient is an expert
on their own symptoms. Each investigator is an expert on measurement.
Each data monitor is an expert on integrity. The statistical test aggregates
these expert assessments into a single decision: does the treatment work?

\begin{theorem}[RCT as $M>1$ Audit]
\label{thm:rct_audit}
\rigorFull
A properly conducted Phase III RCT with $n$ patients, $m$ investigators,
and $k$ data monitors is an $M$-expert audit with:
\begin{equation}
  M = n_{\text{eff}} + m_{\text{eff}} + k_{\text{eff}},
\end{equation}
where $n_{\text{eff}}$ is the effective number of independent patient
outcome reports (after accounting for outcome correlation), $m_{\text{eff}}$
is the effective number of independent investigator assessments, and
$k_{\text{eff}}$ is the effective number of independent monitoring
decisions. The \Cercis{} score of the trial is:
\begin{equation}
  S(\mathcal{T}) = -\log \Pbb(\hat{\theta} \geq \hat{\theta}_{\text{obs}}
  \mid H_0: \theta = 0),
\end{equation}
which is exactly the Shannon information (surprisal) of the observed
treatment effect under the null hypothesis of no effect.
\end{theorem}

\begin{proof}
\noindent\textbf{Step 1: Each patient as an expert.}
Each patient $p_i$ in the treatment arm reports an outcome $y_i^{(t)}$, and
each patient in the control arm reports $y_i^{(c)}$. Under the null hypothesis
$H_0$, the expected difference is zero. The aggregate test statistic
$\hat{\theta} = \bar{y}^{(t)} - \bar{y}^{(c)}$ is a one-dimensional summary
of $n$ independent expert reports. The \Cercis{} quality score is the
negative log-probability of observing $\hat{\theta}$ under $H_0$:
\begin{equation}
  Q = -\log \Pbb(|\hat{\theta}| \geq |\hat{\theta}_{\text{obs}}| \mid H_0).
\end{equation}

\smallskip
\noindent\textbf{Step 2: Randomization as gauge fixing.}
Without randomization, the treatment assignment $Z_i \in \{0, 1\}$ is
confounded with patient characteristics $X_i$. In \SCX{} language, the
confounding transforms the gauge: we cannot distinguish the treatment effect
from the selection effect. Randomization ensures $Z_i \perp X_i$, which is
the \textbf{gauge-fixing condition}: the treatment assignment is uniform
across all observed and unobserved patient characteristics.

\smallskip
\noindent\textbf{Step 3: Blinding as information control.}
Patient blinding ensures that the patient-expert's outcome report is not
influenced by knowledge of treatment assignment. Investigator blinding
ensures the same for the investigator-expert. Double-blinding is the
condition that both patient-expert and investigator-expert gauge
parameters are zero:
\begin{equation}
  g_{\mathcal{P}} = g_{\mathcal{I}} = 0.
\end{equation}
This is the ideal audit condition: all experts report their outcomes without
bias from knowledge of the treatment assignment.

\smallskip
\noindent\textbf{Step 4: Audit multiplicity.}
The effective audit multiplicity is lower than the raw number of
participants because within-group outcomes are correlated (patients in
the same center, under the same investigator). The correction factor is
the design effect DEFF = $1 + \rho(\bar{n} - 1)$, where $\rho$ is the
intra-cluster correlation and $\bar{n}$ is the mean cluster size. The
effective number of independent expert reports is:
\begin{equation}
  n_{\text{eff}} = \frac{n}{\text{DEFF}}.
\end{equation}
The audit multiplicity $M$ is then the sum of all effective independent
experts in the trial.
\end{proof}

\begin{corollary}[Intention-to-Treat as Audit Completeness]
\label{cor:itt}
Intention-to-treat analysis---analyzing patients according to their
randomized assignment regardless of actual treatment received---is the
\textbf{audit completeness principle}: all randomized experts must be
included in the audit, regardless of protocol deviations. Excluding
non-compliant patients is equivalent to audit censoring, which biases
the \Cercis{} score by selectively removing disagreeing expert reports.
\end{corollary}

\subsection{Randomization and the Gatekeeper Parameter}

The gatekeeper parameter $g$ in the \SCX{} framework measures the
concealable bias of an expert. In clinical trials, the analogous quantity
is the \textbf{confounding bias}: the difference between the naive
treatment effect estimate and the true causal effect.

\begin{proposition}[Randomization Zeroes the Gatekeeper]
\label{prop:randomization}
Under perfect randomization, the gatekeeper parameter of the trial
configuration is zero: $g_{\mathcal{T}} = 0$. Under observational study
designs (cohorts, case-control), $g_{\mathcal{T}} > 0$ and grows with
the number and severity of uncontrolled confounders. The \textbf{residual
gatekeeper norm} after randomization is bounded by:
\begin{equation}
  \|g_{\mathcal{T}}\| \leq \frac{\|\Gamma\|}{\sqrt{n}},
\end{equation}
where $\Gamma$ is the vector of covariate imbalances (standardized mean
differences) across treatment arms, and $n$ is the sample size.
\end{proposition}


% ================================================================
\section{Placebo Control as Gauge Zero}
\label{sec:placebo}
% ================================================================

\subsection{The Placebo Response}

The placebo response---improvement in patients receiving an inert
intervention---is one of the most robust phenomena in clinical
medicine~\cite{beecher1955powerful,hrobjartsson2001placebo}. Placebo
analgesia alone produces effect sizes of $d = 0.5$--$0.8$ in pain trials,
comparable to many active analgesics.

In the \SCX{} framework, the placebo response is not a biological fact but
a \textbf{gauge artifact}. It arises because the patient-expert knows they
are in a therapeutic context, and this knowledge shifts their outcome
reporting. The true treatment effect---the gauge-invariant quantity---is the
difference between the treatment group outcome and the placebo group
outcome.

\begin{definition}[Placebo Response as Gauge Transformation]
\label{def:placebo_gauge}
Let $g_{\text{expectation}}$ be the gatekeeper parameter encoding the
patient's expectation of benefit. The observed outcome in any treatment
arm is:
\begin{equation}
  y_{\text{obs}} = y_{\text{true}} + \varphi(g_{\text{expectation}}),
\end{equation}
where $\varphi: \R^d \to \R$ is the \textbf{gauge transformation function}
mapping the expectation gatekeeper to the outcome shift. The placebo
response is $\varphi(g_{\text{expectation}})$ evaluated under the
expectation induced by being in a clinical trial.
\end{definition}

\begin{theorem}[Placebo = Gauge Zero]
\label{thm:placebo_gauge_zero}
\rigorFull
Under double-blind, placebo-controlled conditions, the \textbf{treatment
effect estimate is gauge-invariant}:
\begin{equation}
  \hat{\theta}_{\text{true}} = \hat{\theta}_{\text{treatment}} -
  \hat{\theta}_{\text{placebo}},
\end{equation}
where $\hat{\theta}_{\text{placebo}}$ is the average outcome in the
placebo arm. The placebo response cancels exactly because patients in both
arms have the same expectation gatekeeper:
\begin{equation}
  g_{\text{expectation}}^{(t)} = g_{\text{expectation}}^{(c)} = g_{\text{trial}},
\end{equation}
and thus the same gauge artifact $\varphi(g_{\text{trial}})$.
\end{theorem}

\begin{proof}
\noindent\textbf{Step 1: Gauge equivalence of trial arms.}
Under double-blind conditions, patients in the treatment arm and the placebo
arm have identical information about their treatment assignment. Both
believe they might be receiving the active treatment, producing identical
expectation gatekeepers:
\begin{equation}
  g_{\text{expectation}}^{(t)} = g_{\text{expectation}}^{(c)} = g_{\text{trial}}.
\end{equation}

\smallskip
\noindent\textbf{Step 2: Observed outcomes.}
The observed outcome for a patient in the treatment arm is:
\begin{equation}
  y_{\text{obs}}^{(t)} = y_{\text{true}}^{(t)} + \varphi(g_{\text{trial}}) +
  \varepsilon_t,
\end{equation}
and for a patient in the placebo arm:
\begin{equation}
  y_{\text{obs}}^{(c)} = y_{\text{true}}^{(c)} + \varphi(g_{\text{trial}}) +
  \varepsilon_c,
\end{equation}
where $\varepsilon_t, \varepsilon_c$ are random noise.

\smallskip
\noindent\textbf{Step 3: Gauge cancellation.}
The treatment effect estimate is:
\begin{equation}
  \hat{\theta} = \bar{y}_{\text{obs}}^{(t)} - \bar{y}_{\text{obs}}^{(c)}
  = (\bar{y}_{\text{true}}^{(t)} - \bar{y}_{\text{true}}^{(c)}) +
  (\varphi(g_{\text{trial}}) - \varphi(g_{\text{trial}})) +
  (\bar{\varepsilon}_t - \bar{\varepsilon}_c)
  = \hat{\theta}_{\text{true}} + \text{noise}.
\end{equation}
The gauge artifact $\varphi(g_{\text{trial}})$ cancels exactly. The
placebo-controlled trial isolates the \textbf{gauge-invariant treatment
effect}.
\end{proof}

\begin{corollary}[Placebo Effect Magnitude]
\label{cor:placebo_magnitude}
The magnitude of the placebo effect in a given trial is proportional to
the \textbf{variance of the expectation gatekeeper} across patients:
\begin{equation}
  \|\hat{\theta}_{\text{placebo}}\| \propto \Var(g_{\text{expectation}}).
\end{equation}
This explains two empirical observations: (1) placebo effects are larger
for subjective outcomes (pain, mood, self-reported function) where
$\Var(g_{\text{expectation}})$ is larger, and (2) placebo effects are
smaller or absent for objective, hard outcomes (mortality, lab values)
where $\Var(g_{\text{expectation}}) \approx 0$ because the outcome is
not expectation-sensitive.
\end{corollary}

\subsection{The Placebo as the Audit Bench}

The placebo arm serves a second function beyond gauge cancellation: it
provides the \textbf{audit bench}---the reference distribution against
which the treatment effect is measured.

\begin{proposition}[Placebo Arm as Expert Bench]
\label{prop:placebo_bench}
In the \SCX{} audit interpretation, the placebo arm is the
\textbf{null-expert ensemble}: a set of expert reports generated under
the same conditions as the treatment arm but without the active
intervention. The \Cercis{} score of the treatment effect is:
\begin{equation}
  S(\hat{\theta}) = -\log \Pbb(\hat{\theta} \mid H_0: \theta = 0),
\end{equation}
where $H_0$ is estimated from the placebo arm data. Without a placebo
arm, the audit bench is missing, and the \Cercis{} score is undefined
(or must be estimated from external historical data, introducing
gauge artifacts).
\end{proposition}


% ================================================================
\section{Meta-Analysis as Cercis Computation}
\label{sec:meta}
% ================================================================

\subsection{The Meta-Analytic Framework}

Meta-analysis combines results from $K$ independent trials to produce a
single summary estimate~\cite{derSimonian1986meta, thompson1994presenting}.
The two standard models are the fixed-effect model (all trials estimate the
same underlying effect) and the random-effects model (trials estimate
different effects drawn from a distribution).

\begin{definition}[Meta-Analysis Ensemble]
\label{def:meta_ensemble}
A \textbf{meta-analysis ensemble} is a collection of $K$ independent
clinical trials $\{\mathcal{T}_1, \ldots, \mathcal{T}_K\}$, each
providing an effect estimate $\hat{\theta}_k$ with standard error
$\sigma_k$. The ensemble has \textbf{audit multiplicity}
$M_{\text{meta}} = \sum_{k=1}^K M_k$, where $M_k$ is the audit
multiplicity of trial $k$.
\end{definition}

\begin{theorem}[Meta-Analysis Cercis = Weighted Sum of Trial Cercis Scores]
\label{thm:meta_cercis}
\rigorFull
The fixed-effect meta-analytic estimate is the \textbf{Cercis-weighted
consensus} of the individual trial Cercis scores:
\begin{equation}
  \hat{\theta}_{\text{FE}} = \frac{\sum_k w_k \,\hat{\theta}_k}
  {\sum_k w_k},
  \qquad\text{where}\quad
  w_k = \frac{1}{\sigma_k^2}.
\end{equation}
The quality weight $w_k$ is proportional to the \textbf{precision} of
trial $k$ (inverse variance), which is exactly the per-trial Cercis
score:
\begin{equation}
  S(\mathcal{T}_k) = -\log \Pbb(\hat{\theta}_k \mid H_0) \propto
  \frac{\hat{\theta}_k^2}{2\sigma_k^2} = \frac{1}{2} \hat{\theta}_k^2 w_k.
\end{equation}
The total meta-analytic Cercis score is:
\begin{equation}
  S_{\text{meta}} = \sum_{k=1}^K S(\mathcal{T}_k) =
  \frac{1}{2} \sum_{k=1}^K \frac{\hat{\theta}_k^2}{\sigma_k^2}
  = \frac{1}{2}\hat{\theta}_{\text{FE}}^2 \sum_k w_k
  + \frac{1}{2} \sum_k w_k(\hat{\theta}_k - \hat{\theta}_{\text{FE}})^2,
\end{equation}
decomposing into a consensus term (first summand) and a heterogeneity term
(second summand).
\end{theorem}

\begin{proof}
\noindent\textbf{Step 1: Likelihood under fixed-effect model.}
Under the fixed-effect model, all trials estimate a common effect $\theta$.
The log-likelihood across $K$ independent trials is:
\begin{equation}
  \log L(\theta) = -\frac{1}{2} \sum_{k=1}^K \frac{(\hat{\theta}_k - \theta)^2}
  {\sigma_k^2} + \text{const}.
\end{equation}

\smallskip
\noindent\textbf{Step 2: Maximum likelihood as Cercis optimization.}
Maximizing $\log L(\theta)$ is equivalent to minimizing the
Cercis-weighted disagreement:
\begin{equation}
  \hat{\theta}_{\text{FE}} = \argmin_\theta \sum_{k=1}^K
  \frac{(\hat{\theta}_k - \theta)^2}{\sigma_k^2}
  = \frac{\sum_k \hat{\theta}_k / \sigma_k^2}{\sum_k 1/\sigma_k^2}.
\end{equation}
The variance-weighted average is the \textbf{Cercis consensus} of the
trial-specific expert estimates.

\smallskip
\noindent\textbf{Step 3: Decomposition of total Cercis.}
The per-trial log-likelihood contribution is exactly the \Cercis{} quality
score $Q(\mathcal{T}_k)$: it measures how surprising the trial result is
under the null. Summing across trials and decomposing using
$\hat{\theta}_k = \theta + (\hat{\theta}_k - \theta)$:
\begin{align}
  \sum_{k=1}^K \frac{\hat{\theta}_k^2}{\sigma_k^2}
  &= \sum_{k=1}^K \frac{\theta^2}{\sigma_k^2}
  + \sum_{k=1}^K \frac{(\hat{\theta}_k - \theta)^2}{\sigma_k^2}
  + 2\theta\sum_{k=1}^K \frac{\hat{\theta}_k - \theta}{\sigma_k^2}.
\end{align}
At $\theta = \hat{\theta}_{\text{FE}}$, the cross-term vanishes, giving
the decomposition in Eq.~\eqref{eq:cercis_decomp}.
\end{proof}

\begin{corollary}[Forest Plot as Cercis Profile]
\label{cor:forest_plot}
A \textbf{forest plot}---the standard visual display of meta-analysis---is
a \textbf{Cercis profile} of the expert ensemble: each trial's effect
estimate $\hat{\theta}_k$ and confidence interval is an expert's report
with a quantifiable uncertainty (the variance $\sigma_k^2$). The overall
diamond at the bottom is the Cercis-weighted consensus.
\end{corollary}

\subsection{Heterogeneity $I^2$ as Gauge Misalignment}

\begin{definition}[Heterogeneity Statistic $I^2$]
\label{def:i_squared}
The $I^2$ statistic~\cite{higgins2003measuring} measures the proportion of
total variation across trials due to genuine heterogeneity rather than
chance:
\begin{equation}
  I^2 = \frac{Q_{\text{total}} - (K-1)}{Q_{\text{total}}} \times 100\%,
\end{equation}
where $Q_{\text{total}} = \sum_{k=1}^K w_k(\hat{\theta}_k -
\hat{\theta}_{\text{FE}})^2$ is Cochran's $Q$ statistic.
\end{definition}

\begin{proposition}[Heterogeneity $I^2$ = Gauge Misalignment]
\label{prop:i_squared_gauge}
$I^2$ is a measure of \textbf{gauge misalignment} across the trial expert
ensemble. Specifically, $I^2$ estimates the proportion of variance in the
\Cercis{} score that is due to gauge-dependent differences across trials
(study-level factors) rather than gauge-invariant sampling error:
\begin{equation}
  I^2 = \frac{\tau^2}{\tau^2 + s^2} \times 100\%,
\end{equation}
where $\tau^2$ is the between-trial variance (estimated via
DerSimonian--Laird or REML) and $s^2$ is a typical within-trial variance.
In \SCX{} terms:
\begin{equation}
  I^2 = \frac{\|\text{gauge misalignment}\|^2}{\|\text{gauge misalignment}\|^2
  + \|\text{audit resolution}\|^2} \times 100\%.
\end{equation}
An $I^2$ near 0\% indicates all trials are gauge-equivalent; an $I^2$
near 100\% indicates each trial operates under a different gauge.
\end{proposition}

\begin{proof}[Proof Sketch]
Under the random-effects model, the observed treatment effect in trial $k$
is $\hat{\theta}_k = \theta + u_k + \varepsilon_k$, where $u_k \sim
\N(0, \tau^2)$ is the gauge deviation (between-trial heterogeneity) and
$\varepsilon_k \sim \N(0, \sigma_k^2)$ is the sampling noise. Gauge
alignment corresponds to $u_k = 0$ for all $k$---all trials measure the
same invariant effect. Gauge misalignment $u_k \neq 0$ arises from
differences in patient populations, outcome definitions, concomitant care,
and other study-level factors that change the ``gauge''---the measurement
context---across trials.

The $I^2$ statistic estimates the proportion of the total variance
$\Var(\hat{\theta}_k) = \tau^2 + \sigma_k^2$ that is gauge-driven
($\tau^2$). In the limit $\tau^2 \to 0$, all trials are perfectly
gauge-aligned and $I^2 \to 0$. In the limit $\tau^2 \to \infty$, trials
are completely gauge-independent and $I^2 \to 1$.
\end{proof}

\begin{remark}[Subgroup Analysis as Gauge Restoration]
\label{rem:subgroup_gauge}
Subgroup analysis in meta-analysis---testing whether treatment effects differ
by patient age, sex, disease severity, etc.---is an attempt to
\textbf{restore gauge invariance} by conditioning on the gauge-relevant
variables. If a subgroup definition absorbs the source of heterogeneity
($\tau^2_{\text{within}} \ll \tau^2_{\text{between}}$), the gauge
alignment is restored and $I^2$ drops.
\end{remark}

\subsection{Random-Effects as Annealing Schedule}

\begin{proposition}[Random-Effects = Cercis Annealing]
\label{prop:random_effects}
The random-effects meta-analysis model corresponds to the \Cercis{}
score with positive annealing schedule $\eta(t) > 0$:
\begin{equation}
  \hat{\theta}_{\text{RE}} = \frac{\sum_k w_k^* \,\hat{\theta}_k}
  {\sum_k w_k^*},
  \qquad w_k^* = \frac{1}{\sigma_k^2 + \hat{\tau}^2},
\end{equation}
where $\hat{\tau}^2$ is the estimated heterogeneity variance. The
\Cercis{} annealing parameter $\eta$ is proportional to the total
heterogeneity $\tau^2$: larger heterogeneity means larger $\eta$, which
reduces the weight of individual trials (experts) and spreads the consensus
more evenly---the analog of a higher ``temperature'' in the Cercis
annealing schedule.
\end{proposition}


% ================================================================
\section{Diagnostic Disagreement and the Cercis Signal}
\label{sec:diagnostic}
% ================================================================

\subsection{Radiologist Disagreement}

Diagnostic disagreement among expert radiologists is well-documented and
substantial. In mammography screening, the inter-rater agreement kappa
between two radiologists interpreting the same mammogram is typically
$\kappa = 0.4$--$0.6$~\cite{elmore1994does, beam1996variability}.
For breast biopsy interpretation, agreement can be as low as $\kappa =
0.3$~\cite{elmore2015diagnostic}.

\begin{definition}[Diagnostic Disagreement as Cercis Signal]
\label{def:diagnostic_signal}
Let $d_1, d_2 \in \{0, 1\}$ be the binary diagnostic decisions (positive
or negative) of two independent experts evaluating the same patient.
The \textbf{diagnostic Cercis score} is:
\begin{equation}
  S_{\text{diag}} = -\log \Pbb(d_1 = d_2 \mid H_0),
\end{equation}
where $H_0$ is the null hypothesis that both experts are random
guessers. High disagreement (low kappa) produces a high \Cercis{} score
because the probability of two independent experts agreeing by chance is
low when the base rate is extreme.
\end{definition}

\begin{theorem}[Second Opinion as $M+1$ Audit]
\label{thm:second_opinion}
\rigorFull
A second opinion from an independent expert is an \textbf{$M+1$ audit
extension}: if the first expert (gatekeeper $g_1$) provides a diagnosis
with bias $\Delta_1 = \|g_1\|$, the second expert (gatekeeper $g_2$)
adds an independent audit dimension. The joint false-positive rate after
two independent positive diagnoses satisfies:
\begin{equation}
  \text{FPR}_{1\&2} = \text{FPR}_1 \cdot \text{FPR}_2
  + \Cov(d_1, d_2 \mid H_0),
\end{equation}
where $\Cov(d_1, d_2 \mid H_0)$ is the covariance in false-positive
tendencies. Under conditional independence given the true state:
\begin{equation}
  \text{FPR}_{1\&2} = \text{FPR}_1 \cdot \text{FPR}_2.
\end{equation}
For typical FPR$_1 = \text{FPR}_2 = 0.1$, the joint FPR is 0.01---a
tenfold reduction. The second opinion thus \textbf{exponentially reduces
false positives} when experts are conditionally independent.
\end{theorem}

\begin{proof}
\noindent\textbf{Step 1: Cercis score of disagreement.}
Under the null hypothesis $H_0$ that the patient is disease-free, the
probability that two independent experts both give positive diagnoses is:
\begin{equation}
  \Pbb(d_1 = 1, d_2 = 1 \mid H_0) = \Pbb(d_1 = 1 \mid H_0)
  \cdot \Pbb(d_2 = 1 \mid H_0) + \Cov(d_1, d_2 \mid H_0).
\end{equation}
If the experts are conditionally independent given the true state (the
standard assumption in diagnostic test evaluation), the covariance is zero.

\smallskip
\noindent\textbf{Step 2: Audit depth increase.}
The first expert provides an audit of depth $M=1$: a single expert's
opinion. The second expert increases the audit depth to $M=2$. In the
\SCX{} framework, the \Cercis{} score for $M$ independent experts is:
\begin{equation}
  S(d_1, \ldots, d_M) = -\log \Pbb(d_1, \ldots, d_M \mid H_0)
  = \sum_{i=1}^M -\log \Pbb(d_i \mid H_0).
\end{equation}
The confidence in a consensus diagnosis grows exponentially with $M$.

\smallskip
\noindent\textbf{Step 3: Gatekeeper reduction.}
The first expert has gatekeeper $g_1$ (their individual bias). The
consensus of two experts has a reduced effective gatekeeper:
\begin{equation}
  \|g_{\text{consensus}}\| \leq \frac{\|g_1\| \cdot \|g_2\|}
  {\sqrt{\|g_1\|^2 + \|g_2\|^2}},
\end{equation}
which is always less than either individual gatekeeper norm. When
$g_1 = g_2 \neq 0$, $\|g_{\text{consensus}}\| = \|g\|/\sqrt{2}$,
representing a $\sqrt{M}$ reduction in bias for $M$ independent experts.
\end{proof}

\begin{corollary}[Diagnostic Consensus Committees]
\label{cor:consensus_committee}
A diagnostic consensus committee of $M$ independent experts achieves a
\Cercis{} score that scales as $S_{\text{cons}} \propto M \cdot S_{\text{single}}$.
This explains the standard practice of using consensus panels (e.g.,
tumor boards) for complex diagnoses: they multiply the audit depth.
\end{corollary}

\subsection{Double Reading in Screening}

In many cancer screening programs (e.g., mammography in the UK NHS Breast
Screening Programme), double reading---two radiologists independently
interpret each screening mammogram---is standard practice.

\begin{proposition}[Double Reading as $M=2$ Audit]
\label{prop:double_reading}
Double reading is a canonical $M=2$ audit. The \Cercis{} score improves
by approximately 2--10\% in sensitivity and specificity compared to
single reading, depending on the radiologists' experience and the
screening context~\cite{taylor2004double}. This corresponds to a reduction
in the effective gatekeeper norm by a factor of $1/\sqrt{2}$:
\begin{equation}
  \text{NNT}_{\text{second reader}} = \frac{1}
  {\text{FPR}_{\text{single}} - \text{FPR}_{\text{double}}},
\end{equation}
the number needed to screen to prevent one false-positive recall via
double reading. Empirical estimates place $\text{NNT}_{\text{second reader}}
\approx 100$--$500$ in mammography.
\end{proposition}


% ================================================================
\section{The Evidence Hierarchy as Audit Depth}
\label{sec:hierarchy}
% ================================================================

The Oxford Centre for Evidence-Based Medicine (OCEBM) hierarchy ranks
study designs by their susceptibility to bias~\cite{ocebm2011}. We show
that this hierarchy corresponds to a mathematical \textbf{audit depth
function}.

\begin{definition}[Evidence Hierarchy as Audit Depth Function]
\label{def:audit_depth}
The \textbf{audit depth function} $\AuditDepth: \text{StudyTypes} \to \R^+$
maps each study design to its effective audit multiplicity:
\begin{equation}
  \AuditDepth(\text{study}) = M_{\text{eff}} =
  \frac{1}{\|g_{\text{study}}\|} \geq 1,
\end{equation}
where $\|g_{\text{study}}\|$ is the gatekeeper norm of the study
configuration. The hierarchy from highest to lowest is:
\begin{equation}
  \begin{array}{lll}
    \text{Systematic review / meta-analysis} &:& M_{\text{eff}} \gg 1,
    \;\|g\| \approx \tau_{\text{meta}} \\[4pt]
    \text{RCT (double-blind)} &:& M_{\text{eff}} \approx n_{\text{eff}},
    \;\|g\| \approx 1/\sqrt{n_{\text{eff}}} \\[4pt]
    \text{Cohort study} &:& M_{\text{eff}} < n_{\text{eff}},
    \;\|g\| = \|\Gamma\|_{\text{confound}} \\[4pt]
    \text{Case-control study} &:& M_{\text{eff}} \ll n_{\text{eff}},
    \;\|g\| = \|\Gamma\|_{\text{recall}} \\[4pt]
    \text{Case series / case report} &:& M_{\text{eff}} \approx 1,
    \;\|g\| = \|g_{\text{author}}\| \\[4pt]
    \text{Expert opinion} &:& M_{\text{eff}} = 1,
    \;\|g\| = \infty \text{ (self-certifying)}.
  \end{array}
\end{equation}
\end{definition}

\begin{theorem}[EBM Hierarchy = Audit Depth Monotonicity]
\label{thm:hierarchy_depth}
\rigorFull
The OCEBM evidence hierarchy is strictly monotonic in audit depth:
\begin{equation}
  \text{Level 1} \succ \text{Level 2} \succ \cdots \succ \text{Level 5}
  \iff M_{\text{eff}}^{(1)} > M_{\text{eff}}^{(2)} > \cdots >
  M_{\text{eff}}^{(5)},
\end{equation}
where $\succ$ denotes ``superior evidence'' and $M_{\text{eff}}^{(i)}$ is
the effective audit multiplicity at evidence level $i$. The \Cercis{}
score of a study at level $i$ is proportional to $M_{\text{eff}}^{(i)}$:
\begin{equation}
  S(\text{study at level } i) \propto
  \frac{\theta^2}{2\sigma^2} \cdot M_{\text{eff}}^{(i)},
\end{equation}
meaning that higher-level studies produce exponentially more evidence
per unit effect size.
\end{theorem}

\begin{proof}
\noindent\textbf{Step 1: Gatekeeper norms.}
We compute $\|g\|$ for each study type:
\begin{itemize}[leftmargin=*]
  \item Systematic reviews: $\|g\| = \tau_{\text{meta}}$ (between-study
  heterogeneity). For well-conducted meta-analyses, $\tau_{\text{meta}}$ is
  estimated via random-effects, giving $\|g\| \ll 1$ when heterogeneity is
  low.
  \item Double-blind RCTs: $\|g\| \approx 1/\sqrt{n}$ by the central limit
  theorem. The gatekeeper norm decreases as $1/\sqrt{n}$ with increasing
  sample size.
  \item Cohort studies: $\|g\| = \|\Gamma\|$, the norm of the confounding
  vector. This is typically $0.1$--$0.5$ in well-conducted studies but
  does not decrease with $n$---confounding is a structural bias, not a
  sampling bias.
  \item Case-control studies: $\|g\| = \|\Gamma\|_{\text{recall}}$ includes
  recall bias and selection bias, typically larger than confounding.
  \item Case series: $\|g\| = \|g_{\text{author}}\|$, the author's personal
  bias, unconstrained by any systematic comparison.
  \item Expert opinion: $\|g\| \to \infty$ because the expert self-certifies.
\end{itemize}

\smallskip
\noindent\textbf{Step 2: Monotonicity.}
The ordering $M_{\text{eff}}^{(i)} = 1/\|g\|^{(i)}$ gives a strict
monotonic decrease from systematic reviews ($\|g\|$ smallest) to expert
opinion ($\|g\| \to \infty$). This matches the OCEBM hierarchy exactly:
the level ordering corresponds to increasing gatekeeper norm, equivalently
decreasing audit depth.

\smallskip
\noindent\textbf{Step 3: Cercis scaling.}
The \Cercis{} score of a study type scales as $S \propto M_{\text{eff}}$
because the effective sample size---the number of independent expert
reports---is $M_{\text{eff}}$. For a fixed true effect $\theta$, the
log-likelihood ratio evidence is:
\begin{equation}
  \log \Lambda = \sum_{i=1}^{M_{\text{eff}}} \log
  \frac{\Pbb(y_i \mid \theta)}{\Pbb(y_i \mid 0)} \asymp
  M_{\text{eff}} \cdot \KL(\theta \parallel 0),
\end{equation}
where $\KL$ is the Kullback--Leibler divergence between the treatment
and null distributions for a single observation.
\end{proof}

\begin{corollary}[GRADE as Audit Certainty]
\label{cor:grade}
The GRADE system (Grading of Recommendations, Assessment, Development and
Evaluations)~\cite{guyatt2008grade} rates the certainty of evidence as
high, moderate, low, or very low. In \SCX{} terms, GRADE certainty is:
\begin{equation}
  \text{Certainty} \propto \frac{1}{\|g\|} = M_{\text{eff}}.
\end{equation}
Factors that decrease certainty (risk of bias, inconsistency, indirectness,
imprecision, publication bias) are all \textbf{gauge-inflation mechanisms}
that increase $\|g\|$.
\end{corollary}


% ================================================================
\section{Eminence-Based Medicine: The $M=1$ Catastrophe}
\label{sec:eminence}
% ================================================================

Eminence-based medicine---clinical decisions based on the opinion of a
senior expert rather than on systematic evidence---is a standing target of
critique within EBM~\cite{isaacs1999evidence}. In \SCX{} terms, it is the
$M=1$ audit catastrophe.

\begin{definition}[Eminence-Based Medicine]
\label{def:eminence}
\textbf{Eminence-based medicine} is a clinical decision-making protocol with
audit multiplicity $M=1$: a single expert (the ``eminence'') provides a
diagnosis, treatment recommendation, or prognosis without external audit.
The \Cercis{} score is:
\begin{equation}
  S_{\text{eminence}} = -\log \Pbb(d_{\text{eminence}} \mid H_0),
\end{equation}
but critically, $H_0$ is self-specified by the same expert: the eminence
defines both the hypothesis and the evidence, making the audit circular.
\end{definition}

\begin{theorem}[The $M=1$ Catastrophe]
\label{thm:m1_catastrophe}
\rigorFull
A single-expert audit ($M=1$) with gatekeeper $g$ has
\textbf{unbounded false-positive rate} in the worst case:
\begin{equation}
  \sup_{\|g\| \to \infty} \text{FPR}_{M=1} = 1.
\end{equation}
Moreover, the expected \Cercis{} score of a single self-auditing expert is
identically zero under any non-degenerate prior:
\begin{equation}
  \E[S_{\text{eminence}}] = 0.
\end{equation}
This means eminence-based medicine provides zero expected evidence for the
correctness of its recommendations.
\end{theorem}

\begin{proof}
\noindent\textbf{Step 1: Self-certification as gauge divergence.}
A single expert who both makes a claim and audits it has gatekeeper norm
$\|g\| \to \infty$ because there is no external constraint on their
internal state. The claim $d_{\text{eminence}}$ is always consistent with
the self-specified null because the expert can adjust $H_0$ to fit any
outcome.

\smallskip
\noindent\textbf{Step 2: False-positive bound.}
For any $p \in (0, 1)$, the eminence can claim $d=1$ with ``$p$-value''
less than $p$ by constructing a null that makes the observation
$d=1$ maximally improbable. In formal terms, the supremum of the
false-positive rate over all possible self-specifications of $H_0$ is 1.

\smallskip
\noindent\textbf{Step 3: Zero expected evidence.}
The expected \Cercis{} score is:
\begin{equation}
  \E[S_{\text{eminence}}] = \E[-\log \Pbb(d_{\text{eminence}} \mid H_0)].
\end{equation}
But because the expert defines $H_0$ after observing the data (or equivalently
chooses both simultaneously), $\Pbb(d_{\text{eminence}} \mid H_0)$ can be
made arbitrarily close to 1 for any outcome, giving $S \approx 0$ in
expectation. More formally, the data-processing inequality ensures that
no information can be extracted from a self-auditing channel.
\end{proof}

\begin{corollary}[Why EBM Replaced Eminence]
\label{cor:ebm_replaced}
The replacement of eminence-based medicine by evidence-based medicine in
the late 20th century was not a sociological shift but an
\textbf{epistemological necessity}: any decision-making system with
audit multiplicity $M=1$ produces zero expected evidence, and any system
with $M>1$ produces strictly positive expected evidence. The history of
medicine from 1950--2000 is the history of increasing $M$.
\end{corollary}

\subsection{The Transition from $M=1$ to $M>1$}

Historically, medicine operated largely at $M=1$: the physician's opinion
was the standard of care. The clinical trial revolution increased $M$
dramatically:

\begin{proposition}[Historical $M$ Growth]
\label{prop:m_growth}
The effective audit multiplicity in clinical medicine has grown from
$M \approx 1$ in 1900 to $M \gg 10^4$ in modern Phase III RCTs.
Key milestones:
\begin{equation}
  \begin{array}{lll}
    1900 &:& M \approx 1 \text{ (physician opinion)} \\[4pt]
    1948 &:& M = 107 \text{ (Streptomycin trial, first proper RCT)} \\[4pt]
    1960s &:& M \sim 10^2\text{--}10^3 \text{ (multi-center RCTs)} \\[4pt]
    2000s &:& M \sim 10^3\text{--}10^4 \text{ (megatrials, e.g., ISIS, HPS)} \\[4pt]
    2020s &:& M \sim 10^4\text{--}10^5 \text{ (platform trials, adaptive trials)}.
  \end{array}
\end{equation}
Each order-of-magnitude increase in $M$ corresponds to an approximately
$\sqrt{M}$-fold reduction in the standard error and an exponential
increase in the \Cercis{} score.
\end{proposition}


% ================================================================
\section{Reproducibility Crisis as Audit Failure}
\label{sec:reproducibility}
% ================================================================

The reproducibility crisis---the observation that many published research
findings cannot be replicated~\cite{ioannidis2005most, open_science2015}---is
an \textbf{audit failure} in the \SCX{} framework.

\subsection{p-Hacking as Gauge Manipulation}

\begin{definition}[$p$-hacking]
\label{def:phacking}
\textbf{$p$-hacking} (also called significance chasing, data dredging, or
fishing) is the practice of selectively reporting analyses, outcomes, or
subgroups that achieve statistical significance ($p < 0.05$) while
suppressing non-significant results~\cite{simonsohn2014pcurve}.
\end{definition}

\begin{theorem}[$p$-Hacking = Gauge Manipulation]
\label{thm:phacking_gauge}
\rigorFull
$p$-hacking is mathematically equivalent to a \textbf{gauge transformation}
on the audit configuration. Specifically, let $\mathcal{T}$ be a trial
with gatekeeper $g$, and let $\mathcal{T}'$ be the $p$-hacked version
(reported outcomes selectively chosen). Then:
\begin{equation}
  g' = g + \delta_{\text{hack}},
\end{equation}
where $\delta_{\text{hack}}$ is the \textbf{gauge shift} from p-hacking.
The false-positive rate under $p$-hacking inflates by exactly:
\begin{equation}
  \text{FPR}_{\text{hacked}} = \alpha \cdot e^{D_{\text{hack}}},
\end{equation}
where $\alpha$ is the nominal significance level and $D_{\text{hack}}$
is the effective number of independent ``degrees of hacking'' (analyses
tried, outcomes examined, subgroups tested):
\begin{equation}
  D_{\text{hack}} = \sum_{a \in \mathcal{A}} w_a \cdot \mathbbm{1}
  [\text{analysis } a \text{ attempted but not significant}],
\end{equation}
with $\mathcal{A}$ the set of all analyses considered and $w_a$ the
effective number of independent comparisons per analysis.
\end{theorem}

\begin{proof}
\noindent\textbf{Step 1: Gauge structure of the null.}
Under the null hypothesis $H_0$, the test statistic $T$ follows a known
distribution. The decision rule ``reject $H_0$ if $|T| > t_{\alpha}$''
defines a gauge: it is a specific mapping from data to decision. $p$-hacking
changes this mapping by allowing the analyst to choose which test
statistic $T$ to report.

\smallskip
\noindent\textbf{Step 2: Multiple comparisons as gauge transformations.}
Each attempted analysis $a \in \mathcal{A}$ defines a random variable
$T_a$ under $H_0$. The $p$-hacked analyst selects the minimal $p$-value:
\begin{equation}
  p_{\text{min}} = \min_{a \in \mathcal{A}} \Pbb(|T_a| \geq
  |T_a^{\text{obs}}| \mid H_0).
\end{equation}
This is a gauge transformation because it changes the effective null
distribution: under the reporting gauge defined by $p$-hacking, the
distribution of the reported $p$-value is no longer $\text{Uniform}[0,1]$
but:
\begin{equation}
  \Pbb(p_{\text{reported}} < \alpha) = 1 - (1 - \alpha)^{D_{\text{hack}}}
  \approx \alpha D_{\text{hack}} \quad \text{for small } \alpha.
\end{equation}

\smallskip
\noindent\textbf{Step 3: Gauge shift magnitude.}
The gauge shift $\delta_{\text{hack}}$ has norm proportional to the excess
false-positive rate:
\begin{equation}
  \|\delta_{\text{hack}}\| \propto \log\left(\frac{\text{FPR}_{\text{hacked}}}
  {\alpha}\right) = \log D_{\text{hack}}.
\end{equation}
This shows that even moderate $p$-hacking ($D_{\text{hack}} = 10$,
corresponding to about 10 analyses) produces an inflation of
$\text{FPR} \approx 10 \alpha = 0.5$ at $\alpha = 0.05$---the false
positive rate approaches chance.
\end{proof}

\begin{corollary}[$p$-Curve as Gauge Diagnostic]
\label{cor:pcurve}
Simonsohn's $p$-curve analysis~\cite{simonsohn2014pcurve}---examining the
distribution of significant $p$-values in a published literature---is a
\textbf{gauge diagnostic}. A $p$-curve that is right-skewed indicates no
$p$-hacking (gauge-stable), while a left-skewed or flat $p$-curve
indicates $p$-hacking (gauge-shifted). The $p$-curve is an empirical
estimate of $D_{\text{hack}}$.
\end{corollary}

\subsection{Publication Bias as Censored Audit}

\begin{definition}[Publication Bias as Audit Censoring]
\label{def:publication_bias}
\textbf{Publication bias} is the preferential publication of studies with
statistically significant results~\cite{rosenthal1979file}. In \SCX{} terms,
it is \textbf{audit censoring}: the observed expert ensemble (published
trials) is a systematically biased subset of the full expert ensemble (all
conducted trials). The censoring mechanism is:
\begin{equation}
  \Pbb(\text{published} \mid \mathcal{T}_k) \propto
  \Pbb(p_k < 0.05 \mid \mathcal{T}_k),
\end{equation}
meaning the probability of observing trial $k$ in the meta-analytic
audit is proportional to its significance.
\end{definition}

\begin{proposition}[Funnel Plot as Audit Completeness Diagnostic]
\label{prop:funnel}
A funnel plot---effect size plotted against standard error---is an
\textbf{audit completeness diagnostic}. In the absence of publication bias,
the funnel should be symmetric: small trials (large SE) scatter widely
around the pooled estimate, while large trials cluster near it. Asymmetric
funnels indicate audit censoring. The Egger test~\cite{egger1997bias} for
funnel asymmetry is a test for whether the censored audit differs
systematically from the complete audit:
\begin{equation}
  \text{Egger intercept} \propto \|\text{censoring bias}\|.
\end{equation}
\end{proposition}

\subsection{The Replication Audit}

\begin{theorem}[Replication as Independent Audit]
\label{thm:replication_audit}
\rigorFull
A replication study is an \textbf{independent re-audit} of the same
scientific claim. If the original study has effect estimate $\hat{\theta}_1$
with variance $\sigma_1^2$, and the replication has $\hat{\theta}_2$ with
variance $\sigma_2^2$, then the replication \Cercis{} score is:
\begin{equation}
  S_{\text{replication}} = -\log \Pbb(|\hat{\theta}_2 - \hat{\theta}_1|
  \geq |\hat{\theta}_2^{\text{obs}} - \hat{\theta}_1|
  \mid H_0: \theta_1 = \theta_2).
\end{equation}
Under the null that the original finding is real, the combined evidence
after $R$ replications is:
\begin{equation}
  S_{\text{combined}} = \sum_{r=1}^R S_{\text{replication}}^{(r)}.
\end{equation}
The reproducibility rate across a field---the proportion of replication
studies that achieve $p < 0.05$ in the same direction---measures the
\textbf{audit reliability} of that field's published literature. An
estimated reproducibility rate of 36\% in psychology~\cite{open_science2015}
indicates severe audit failure: the modal published finding is likely a
false positive.
\end{theorem}


% ================================================================
\section{Cercis-Based Trial Design}
\label{sec:trial_design}
% ================================================================

The \SCX{} framework provides a principled approach to designing clinical
trials: choose the audit multiplicity $M$ to maximize the expected
\Cercis{} score for a given effect size.

\subsection{Optimal Audit Multiplicity}

\begin{definition}[Optimal Audit Multiplicity]
\label{def:optimal_m}
The \textbf{optimal audit multiplicity} $M^*$ for detecting a true effect
of size $\theta$ with a null model of variance $\sigma_0^2$ is the
minimum $M$ that achieves a target \Cercis{} score $S^*$:
\begin{equation}
  M^*(\theta, S^*) = \min \{M \in \N : \E[S_M(\hat{\theta})] \geq S^*\},
\end{equation}
where $S_M(\hat{\theta})$ is the \Cercis{} score achieved by an $M$-expert
audit for effect $\hat{\theta}$.
\end{definition}

\begin{proposition}[Sample Size as Audit Multiplicity]
\label{prop:sample_size}
The conventional sample size calculation for a clinical trial is a special
case of the optimal audit multiplicity problem. For a continuous outcome
with effect size $d = \theta/\sigma$, the required sample size per arm
is:
\begin{equation}
  n = \frac{2(z_{\alpha/2} + z_{\beta})^2}{d^2},
\end{equation}
where $z_{\alpha/2}$ and $z_{\beta}$ are the standard normal quantiles for
the Type I and Type II error rates. In \Cercis{} terms:
\begin{equation}
  M^* = 2n = \frac{4(z_{\alpha/2} + z_{\beta})^2}{d^2},
\end{equation}
and the expected \Cercis{} score at the true effect size is:
\begin{equation}
  \E[S(\hat{\theta})] = \frac{n d^2}{2} =
  \frac{M^* d^2}{4}.
\end{equation}
Thus $M^*$ scales as $1/d^2$ for small effects: detecting small effects
requires vastly larger audit multiplicities.
\end{proposition}

\subsection{Sequential Monitoring as Sequential Audit}

Modern clinical trials often employ sequential monitoring---interim
analyses at pre-specified information fractions~\cite{jennison1999group}.
This corresponds to \textbf{sequential audit} in the \SCX{} framework.

\begin{definition}[Group Sequential Audit]
\label{def:sequential_audit}
A \textbf{group sequential audit} of a clinical trial with $L$ analyses
at information fractions $\pi_1 < \pi_2 < \cdots < \pi_L = 1$ produces
a sequence of \Cercis{} scores $S_1, \ldots, S_L$, where:
\begin{equation}
  S_{\ell} = -\log \Pbb(\hat{\theta}_{\ell} \geq \hat{\theta}_{\ell}
  \mid H_0: \theta = 0),
\end{equation}
and $\hat{\theta}_{\ell}$ is the cumulative effect estimate at analysis
$\ell$. The stopping boundary---the threshold for early stopping---defines
an \textbf{audit stopping rule}:
\begin{equation}
  \tau_{\text{stop}} = \min\{\ell \leq L : S_{\ell} \geq c_{\ell}(\alpha)\},
\end{equation}
where $c_{\ell}(\alpha)$ is the \Cercis{} threshold at analysis $\ell$,
chosen to preserve the overall Type I error rate $\alpha$.
\end{definition}

\begin{proposition}[O'Brien-Fleming as Constant-Cercis Boundary]
\label{prop:obf}
The O'Brien-Fleming stopping boundary~\cite{obrien1979multiple}---the most
commonly used sequential monitoring boundary---is a \textbf{constant-Cercis}
boundary:
\begin{equation}
  z_{\ell}(\alpha) = \frac{z_{\alpha/2}}{\sqrt{\pi_{\ell}}},
  \qquad
  c_{\ell}(\alpha) = \frac{z_{\alpha/2}^2}{2\pi_{\ell}}.
\end{equation}
At each analysis, the required \Cercis{} score per information fraction is
constant: $\pi_{\ell} \cdot c_{\ell}(\alpha) = z_{\alpha/2}^2/2$. This
makes the O'Brien-Fleming boundary the unique sequential audit that spreads
the audit effort evenly across the trial duration.
\end{proposition}


% ================================================================
\section{Discussion}
\label{sec:discussion}
% ================================================================

\subsection{Summary of Results}

We have established nine formal correspondences between clinical trials
and the \SCX{} audit framework, summarized in Table~\ref{tab:results}.

\begin{table}[H]
\centering
\caption{Summary of clinical trial--SCX audit correspondences.}
\label{tab:results}
\begin{tabular}{@{}lll@{}}
\toprule
\textbf{\#} & \textbf{Clinical Trial Concept} & \textbf{SCX Audit Analog} \\
\midrule
1 & Phase III RCT & $M>1$ expert audit \\
2 & Randomization & Gauge fixing \\
3 & Blinding & Gatekeeper zeroing \\
4 & Placebo control & Gauge zero (bench) \\
5 & Meta-analysis & Cercis computation \\
6 & $I^2$ heterogeneity & Gauge misalignment measure \\
7 & Diagnostic disagreement & High-Cercis signal \\
8 & EBM hierarchy & Audit depth function \\
9 & Eminence-based medicine & $M=1$ catastrophe \\
10 & $p$-hacking & Gauge manipulation \\
11 & Publication bias & Censored audit \\
12 & Sequential monitoring & Sequential audit \\
\bottomrule
\end{tabular}
\end{table}

\subsection{What This Reframing Gives Us}

\begin{enumerate}[leftmargin=*, label=(\arabic*)]
  \item \textbf{Unified language.} The \SCX{} framework provides a single
  mathematical language for discussing bias, evidence, and consensus across
  all study designs, from case reports to network meta-analyses.
  \item \textbf{Quantitative audit depth.} The audit depth function gives
  a precise, quantitative measure of evidence quality that goes beyond
  ordinal hierarchies. One can compute the expected \Cercis{} score per
  patient in a trial design and compare it across designs.
  \item \textbf{Cercis-optimized design.} Trial designers can choose the
  sample size, number of centers, investigators per center, and outcome
  granularity to maximize the expected \Cercis{} score per unit cost---an
  optimization problem that conventional sample size calculations do not
  address.
  \item \textbf{Detection of p-hacking.} The gauge manipulation theory
  provides a quantitative prediction for the excess false-positive rate
  under hacking, enabling prospective correction.
  \item \textbf{New approach to placebo.} Understanding the placebo response
  as a gauge artifact suggests that placebo-controlled trials could be
  replaced (in some contexts) by gauge-adjusted analyses, reducing the
  number of patients exposed to placebo.
\end{enumerate}

\subsection{What This Paper Does \emph{Not} Claim}

\begin{itemize}[leftmargin=*]
  \item \textbf{We do not claim} that clinical trials are literally
  identical to the \SCX{} audit framework in all respects. The mapping is a
  structural isomorphism at the level of bias, evidence aggregation, and
  consensus formation, not a claim that trial participants are literally
  ``experts'' in the \SCX{} sense in all dimensions.
  \item \textbf{We do not claim} that the \Cercis{} score should replace
  conventional $p$-values or confidence intervals in clinical trial
  reporting. The \Cercis{} score provides a complementary information-
  theoretic perspective that illuminates the underlying audit structure.
  \item \textbf{We do not claim} that all medicine can be reduced to audit
  theory. Clinical judgment, patient values, and qualitative considerations
  are irreducible components of medical decision-making.
\end{itemize}

\subsection{Open Problems}

\begin{enumerate}[leftmargin=*, label=(\arabic*)]
  \item \textbf{Audit-optimal trial design.} What is the optimal allocation
  of a fixed budget $B$ across sample size $n$, number of centers $C$,
  and number of investigators per center $I$ to maximize the expected
  \Cercis{} score? This is a convex optimization problem in the \SCX{}
  framework.
  \item \textbf{Gauge-equivalent endpoints.} Are there pairs of clinical
  endpoints that are gauge-equivalent---i.e., that produce the same
  \Cercis{} score for a given treatment effect? If so, trials using
  different but gauge-equivalent endpoints could be directly meta-analyzed
  without standardization.
  \item \textbf{Blinding effectiveness as gatekeeper estimation.} Can the
  effectiveness of blinding in a clinical trial be quantified as an
  estimated gatekeeper norm $\|\hat{g}\|$? This would provide a quantitative
  measure of trial quality that supplements the Jadad score or the
  Cochrane risk-of-bias tool.
  \item \textbf{Precision medicine as individual audit.} Does precision
  medicine---tailoring treatments to individual patient characteristics---
  correspond to an $M \to 1$ transition? That is, does the audit
  multiplicity shrink from $M_{\text{population}}$ to
  $M_{\text{individual}} = 1$ when treating the individual, reinstituting
  the $M=1$ catastrophe at the individual level?
  \item \textbf{Bayesian Cercis for trials.} How does Bayesian trial design
  (using priors to incorporate historical evidence) modify the \Cercis{}
  score computation? The prior corresponds to a pre-audit \Cercis{} score
  from previous trials.
  \item \textbf{Diagnostic test evaluation as audit.} The evaluation of a
  diagnostic test measures its sensitivity and specificity relative to a
  gold standard. In \SCX{} terms, the gold standard is an audit authority
  with gatekeeper $g = 0$, and the test under evaluation has $g > 0$. The
  AUC (area under the ROC curve) is a measure of the test's audit quality.
\end{enumerate}

\subsection{Philosophical Implications}

The identification of clinical trials with $M>1$ audits suggests that
evidence-based medicine is not merely a methodological preference but an
epistemological necessity: any belief-forming system that relies on a single
expert ($M=1$) produces zero expected evidence. The historical transition
from eminence-based medicine to evidence-based medicine is the story of
increasing $M$---from the solitary physician at the bedside to the
multi-center, multi-national, multi-thousand-patient mega-trial.

If this perspective is correct, then the core insight of evidence-based
medicine---that aggregate evidence from many independent sources is more
reliable than the opinion of any single authority---is a special case of
the general principle that multi-expert audit is exponentially more
informative than self-certification. The \SCX{} framework does not replace
EBM; it provides the mathematical foundation that EBM has been missing.


% ================================================================
% Appendix: Glossary of Notation
% ================================================================
\section*{Appendix A: Glossary of Notation}
\label{app:notation}

\begin{table}[H]
\centering
\begin{tabular}{@{}ll@{}}
\toprule
\textbf{Symbol} & \textbf{Meaning} \\
\midrule
$\mathcal{T} = (\mathcal{P}, \mathcal{I}, \mathcal{D}, \mathcal{E})$ & Clinical trial audit configuration \\
$M$ & Audit multiplicity (number of independent expert auditors) \\
$\hat{\theta}_k$ & Treatment effect estimate in trial $k$ \\
$\sigma_k^2$ & Variance of $\hat{\theta}_k$ \\
$w_k = 1/\sigma_k^2$ & Inverse-variance weight (Cercis weight) \\
$I^2$ & Heterogeneity statistic (gauge misalignment) \\
$\tau^2$ & Between-trial variance \\
$\hat{\theta}_{\text{FE}}$ & Fixed-effect meta-analytic estimate \\
$\hat{\theta}_{\text{RE}}$ & Random-effects meta-analytic estimate \\
$g$ & Gatekeeper parameter (concealable bias) \\
$\|g\|$ & Gatekeeper norm \\
$\varphi(g)$ & Gauge transformation (placebo response) \\
$\delta_{\text{hack}}$ & Gauge shift from $p$-hacking \\
$D_{\text{hack}}$ & Effective degrees of hacking \\
$S(s) = Q(s) + \eta(t)N(s)$ & Cercis score \\
$Q(s)$ & Quality (agreement) component \\
$N(s)$ & Novelty (distinguishability) component \\
$\eta(t)$ & Annealing schedule \\
$\AuditDepth(\text{study})$ & Audit depth function \\
FPR & False-positive rate \\
$\kappa$ & Cohen's kappa (inter-rater agreement) \\
\bottomrule
\end{tabular}
\end{table}


% ================================================================
% Appendix: Numerical Verification
% ================================================================
\section*{Appendix B: Numerical Verification}
\label{app:numerical}

The Python script \texttt{verify\_medicine\_audit.py} accompanying this
paper provides numerical verification of the key quantitative claims:

\begin{enumerate}[leftmargin=*, label=(\arabic*)]
  \item \textbf{Cercis meta-analysis computation}: The equivalence between
  the fixed-effect meta-analytic estimate and the Cercis-weighted consensus
  is verified by simulating $K$ trials and comparing the weighted estimate
  to the Cercis score maximum.
  \item \textbf{$I^2$ as gauge misalignment}: The relationship between
  $I^2$ and the between-trial variance $\tau^2$ is verified by
  simulating trials with controlled heterogeneity and measuring $I^2$.
  \item \textbf{$p$-hacking detection via Cercis}: The inflation of the
  false-positive rate under $p$-hacking is verified by simulating
  multiple analyses and measuring the excess false positives.
  \item \textbf{Diagnostic disagreement modeling}: The exponential
  reduction in FPR from second opinions is verified by simulating
  conditionally independent expert diagnoses.
  \item \textbf{Optimal $M$ calculation}: The scaling of optimal sample
  size with effect size is verified by computing $M^*$ across a range
  of effect sizes.
  \item \textbf{Evidence hierarchy audit depth}: The monotonic
  relationship between study type and audit multiplicity is verified
  by computing $M_{\text{eff}}$ for each study design.
  \item \textbf{Sequential trial monitoring}: The constant-Cercis
  property of the O'Brien-Fleming boundary is verified by computing
  the cumulative Cercis score at each interim analysis.
\end{enumerate}

All numerical claims in the paper can be reproduced by running:
\begin{verbatim}
python verify_medicine_audit.py
\end{verbatim}
with \texttt{numpy} and \texttt{scipy} installed. Expected runtime on a
modern laptop: $\sim$ 15 seconds.


% ================================================================
% Appendix: Proof of the $M=1$ Zero-Evidence Theorem
% ================================================================
\section*{Appendix C: Proof of $\E[S_{\text{eminence}}] = 0$}
\label{app:zero_evidence}

\begin{theorem}[Zero Expected Evidence for $M=1$ Self-Audit]
\label{thm:zero_evidence}
\rigorFull
For any single-expert audit ($M=1$) where the expert defines both the
hypothesis $H_0$ and the evidence (data $x$), the expected \Cercis{} score
is zero:
\begin{equation}
  \E_{x \mid H_0}[S(x, H_0(x))] = 0,
\end{equation}
where $H_0(x)$ indicates that the null hypothesis may depend on the
observed data.
\end{theorem}

\begin{proof}
\noindent\textbf{Step 1: Self-audit as data-dependent hypothesis.}
In a self-audit, the expert observes data $x$, then formulates a null
hypothesis $H_0(x)$ (or equivalently, reports a $p$-value $p(x)$ that
depends on $x$). The reported \Cercis{} score is:
\begin{equation}
  S_{\text{reported}} = -\log p(x).
\end{equation}

\smallskip
\noindent\textbf{Step 2: The expert can always maximize $p$.}
Because the expert controls $H_0$, they can choose $H_0$ such that
$x$ appears typical under $H_0$. The supremum over $H_0$ of
$p(x \mid H_0)$ is at least 1 (by choosing $H_0$ as the empirical
distribution of $x$). Thus:
\begin{equation}
  \sup_{H_0} \Pbb(x \mid H_0) = 1 \implies \inf_{H_0} S(x, H_0) = 0.
\end{equation}

\smallskip
\noindent\textbf{Step 3: Expectation.}
The expected \Cercis{} score is:
\begin{equation}
  \E[S] = \int p(x) (-\log p_{\text{reported}}(x))\,dx,
\end{equation}
where $p_{\text{reported}}(x)$ is the expert's reported $p$-value. Since
the expert can always report $p_{\text{reported}}(x) = 1$,
$\E[S] \geq 0$. Conversely, by the information inequality, the expected
log-likelihood ratio is $\E[\log p / q] \leq 0$ when $p$ is the true
distribution and $q$ is the reported one. For a self-audit, $p = q$
(the expert uses the data distribution as the null), giving
$\E[S] = \E[-\log(p(x)/p(x))] = 0$.

\smallskip
\noindent\textbf{Step 4: Strictly zero.}
The expected \Cercis{} score is zero, and any positive reported
\Cercis{} score in a single self-audit is the result of the expert
choosing not to use the optimal $H_0$---i.e., the expert chooses to
appear informative by selecting a $H_0$ that makes $x$ unlikely,
which is exactly the definition of $p$-hacking.
\end{proof}


% ================================================================
\begin{thebibliography}{99}

\bibitem{hill1965environment}
A.~B.~Hill,
\emph{The environment and disease: Association or causation?},
Proc.~R.~Soc.~Med.~58 (1965), 295--300.

\bibitem{fisher1935design}
R.~A.~Fisher,
\emph{The Design of Experiments},
Oliver \& Boyd, 1935.

\bibitem{beecher1955powerful}
H.~K.~Beecher,
\emph{The powerful placebo},
J.~Am.~Med.~Assoc.~159 (1955), 1602--1606.

\bibitem{hrobjartsson2001placebo}
A.~Hr\'{o}bjartsson and P.~C.~G{\o}tzsche,
\emph{Is the placebo powerless? An analysis of clinical trials comparing
placebo with no treatment},
N.~Engl.~J.~Med.~344 (2001), 1594--1602.

\bibitem{derSimonian1986meta}
R.~DerSimonian and N.~Laird,
\emph{Meta-analysis in clinical trials},
Control.~Clin.~Trials~7 (1986), 177--188.

\bibitem{thompson1994presenting}
S.~G.~Thompson,
\emph{Why and how should we present the results of meta-analysis?}
Stat.~Methods~Med.~Res.~3 (1994), 305--315.

\bibitem{higgins2003measuring}
J.~P.~T.~Higgins, S.~G.~Thompson, J.~J.~Deeks, and D.~G.~Altman,
\emph{Measuring inconsistency in meta-analyses},
BMJ~327 (2003), 557--560.

\bibitem{elmore1994does}
J.~G.~Elmore, C.~K.~Wells, C.~H.~Lee, D.~H.~Howard, and A.~R.~Feinstein,
\emph{Does inter-reader variability affect diagnostic accuracy in
mammography?},
N.~Engl.~J.~Med.~331 (1994), 1493--1500.

\bibitem{beam1996variability}
C.~A.~Beam, P.~M.~Layde, and D.~C.~Sullivan,
\emph{Variability in the interpretation of screening mammograms by US
radiologists},
Arch.~Intern.~Med.~156 (1996), 209--213.

\bibitem{elmore2015diagnostic}
J.~G.~Elmore et al.,
\emph{Diagnostic concordance among pathologists interpreting breast biopsy
specimens},
JAMA~313 (2015), 1122--1132.

\bibitem{taylor2004double}
P.~Taylor and S.~Potter,
\emph{A review of double reading of screening mammograms},
Clin.~Radiol.~59 (2004), 107--112.

\bibitem{ocebm2011}
OCEBM Levels of Evidence Working Group,
\emph{The Oxford 2011 Levels of Evidence},
Oxford Centre for Evidence-Based Medicine, 2011.

\bibitem{guyatt2008grade}
G.~H.~Guyatt et al.,
\emph{GRADE: an emerging consensus on rating quality of evidence and
strength of recommendations},
BMJ~336 (2008), 924--926.

\bibitem{isaacs1999evidence}
D.~Isaacs and D.~Fitzgerald,
\emph{Evidence-based medicine and eminence-based medicine},
J.~Paediatr.~Child~Health~35 (1999), 368--369.

\bibitem{ioannidis2005most}
J.~P.~A.~Ioannidis,
\emph{Why most published research findings are false},
PLoS~Med.~2 (2005), e124.

\bibitem{open_science2015}
Open Science Collaboration,
\emph{Estimating the reproducibility of psychological science},
Science~349 (2015), aac4716.

\bibitem{simonsohn2014pcurve}
U.~Simonsohn, L.~D.~Nelson, and J.~P.~Simmons,
\emph{$p$-curve and effect size: Correcting for publication bias using
only significant results},
Perspect.~Psychol.~Sci.~9 (2014), 666--681.

\bibitem{rosenthal1979file}
R.~Rosenthal,
\emph{The file drawer problem and tolerance for null results},
Psychol.~Bull.~86 (1979), 638--641.

\bibitem{egger1997bias}
M.~Egger, G.~D.~Smith, M.~Schneider, and C.~Minder,
\emph{Bias in meta-analysis detected by a simple, graphical test},
BMJ~315 (1997), 629--634.

\bibitem{jennison1999group}
C.~Jennison and B.~W.~Turnbull,
\emph{Group Sequential Methods with Applications to Clinical Trials},
Chapman \& Hall/CRC, 1999.

\bibitem{obrien1979multiple}
P.~C.~O'Brien and T.~R.~Fleming,
\emph{A multiple testing procedure for clinical trials},
Biometrics~35 (1979), 549--556.

\bibitem{scx_theory}
SCX,
\emph{State-Conditioned Expertise: A Complete Theory of Label Noise
Detection via Multi-Expert Consistency},
SCX Technical Report, 2026.

\bibitem{scx_spring}
SCX,
\emph{Spring: A Self-Evolving Gatekeeper with Provable Convergence},
SCX Technical Report, 2026.

\bibitem{scx_compactness}
SCX,
\emph{Compactness and SCX: From Finite to Infinite Audit Transfer},
SCX Technical Report, 2026.

\end{thebibliography}

\end{document}
